\section{Locations of the companion proofs}\label{app:companions}

The moment-defect and original count arguments are due to
Rodziewicz--Ko\l{}odziejek \cite{RK2026}; the completion notes identify
their reconstruction scope. The companion sources retain their original
attribution, theorem labels and statement scopes. A label quoted below belongs to the named companion,
not to the numbering of this article. The source collection is prepared
by \path{prepare_global_companions.py}; it includes the local TeX
inputs of the multifile large-index proof.

\paragraph{Input (I).}
In \path{companions/low-index/}, the first-index bound is
\path{n1/n1_base.tex}, Theorem~1 (\texttt{thm:base}).
The source \path{n2-n3/manuscript.tex} contains
\texttt{n2:thm:main} and \texttt{n3:thm:main}, together with their
support reductions and continuous-domain polynomial certificate
implications. The complete fourth-index proof is available in the
existing public manuscript PDF
\path{companions/references/rk-v3/moment-defects-v3.pdf},
Section~26, Theorem~26.1 and its proof. A TeX source for that complete
fourth-index proof is not included.

The subsequent compositions are Theorem~1
(\texttt{composition}) in \path{n5/five_sum_completion_20261025.tex};
Proposition~1 and its completed composition in each of
\path{n6/six_sum_interface_20261025.tex},
\path{n7/seven_sum_interface_20261025.tex},
\path{n8/eight_sum_interface_20261025.tex}, and
\path{n9/nine_sum_interface_20261025.tex};
the theorem entitled \emph{Unrestricted tenth-index conclusion}
in \path{n10/main.tex}; and the ordered implication at the end of
\path{n11-14/conditional_11_14_interface_20261026.tex}
(which has no separately numbered conclusion).
The final blocks are Theorem~1 in \path{n15-19/main.tex},
Proposition~1 and its completed composition in
\path{n20/twenty_sum_interface_20261025.tex},
Theorem~1 in \path{n21-24/main.tex}, and
Theorem~1 in \path{n25/main.tex}.
The high-moment complement at the actual indices $25$--$30$
is stated and derived in \path{n25-fourier/main.tex}.

These interfaces are accompanied by continuous derivations.
The source
\path{companions/common/finite-analytic/BE_cutoff35_streamlined.tex}
contains \path{global:attainment}, \texttt{global:contact},
\path{complete44ent:smooth-fourier},
\path{complete44ent:fourier-early-low},
\path{complete44ent:re-enclosure}, \texttt{res:contract},
and the section \emph{One finite-count comparison for all the new indices}.
The additional small-index proofs retain their original text format under
\path{companions/low-index/continuous-sources/work/}.
In particular, the fifth-index
\path{cutoff5_20261020/audit/ANALYTIC_DEPENDENCIES_V2.txt}
locates the conserved-identity, finite-window LOW, resource and count
derivations; their mathematical proof texts are included alongside that
source map. They are distinct from numerical success records.
The original frozen arithmetic inputs and actual returned objects are
retained under \path{companions/low-index-support/}.
Its \path{read_low_index_support.py} describes the existing joins,
their exact input locations and their Python or compiler requirements.

\paragraph{Input (II).}
In \path{companions/actual-index/}, the coarse entry from
$C_{30}<\CE+.01$ and the literal bridges to $34$ are
Propositions \texttt{prop:31} and \texttt{prop:bridges} in
\path{n31-34/main.tex}. The exceptional band-$7$ defect-domain
argument is \path{n31-domain/sharp31_B7_D_domain_20261025.tex}.
Continue with the theorem entitled
\emph{Unrestricted index--35 conclusion} in
\path{n35/actual35_conditional_composition_20261025.tex},
the proposition entitled \emph{Unrestricted conclusions at 36, 37 and 38}
in \path{n36-38/actual36_38_conditional_composition_20261026.tex},
and Propositions~1 and~2 in \path{n39/main.tex}.
The actual-$39$ continuous-domain extension is separately included in
\path{n39-continuous/actual39_continuous_lemmas_20261026.tex}.
Theorem \texttt{thm:main} in \path{n40-48/main.tex} contains the
bridge through $48$; the present composition uses its steps only
through $46$. Its two local implications are
\texttt{prop:forty} and \texttt{prop:two}.
The directories \path{n40-continuous/}, \path{n40-original/},
\path{n42-continuous/}, and \path{n42-local/} contain the corresponding
actual-index and finite-band continuous arguments.
The corresponding compressed inputs, returned objects and existing
join programs are retained under \path{companions/bridge-support/}.
The accompanying \path{read_bridge_support.py} locates the mathematical
objects and launches the original joins against those collected inputs.

\paragraph{Input (III).}
The local resource implication \path{re:localG} in
\path{paper_entry/resource_entry.tex} is used with $b\ge0$, $P\ge0$,
$\ell_-\ge0$ and $L_0\ge0$. These are the physical coefficients of
\path{conf:moment} and the nonnegative noise subintervals; they license
the substitutions $nP\ge NP$ and $n\theta\ge\ell_-^2L_0$.

The proposition \texttt{global:local} in
\path{companions/large-index/analytic-route/paper/sections/global_route.tex}
contains the zero-allowance uniform local exclusion (III) used here.
In the present proof, Corollary~\ref{gc:tail-support} replaces its
initial radius-$3.84$ calculation over six index bands, both left radial
contractions and the central radius-$2.5$ contraction. The resource
schedule, coordinate-strip exclusions and local comparisons are
retained. The broader original heat-expansion route remains in the archive.
The continuous analytic route is the multifile source
\path{companions/large-index/analytic-route/paper/main.tex}:
\texttt{re:acquired}, \texttt{re:enclosure}, and
\path{shortloc:theorem} supply its resource acquisition and
local comparison. Its included coefficient atlas derives
compensated inversion and the conserved identity, rather than treating
them as numerical premises. The independently completed $n\ge49$
route is also retained in
\path{companions/large-index/completion/n49_completion_20261025.tex};
the new induction from $45,46$ uses the local exclusion directly.
The selected compressed source/input archive and its original returned
calculation archive are included under \path{companions/uniform-tail/}.
The accompanying \path{prepare_uniform_tail_support.py} can unpack
them to a new directory, reconstruct the exact local cover, and run
the original transcript-to-target join. Its Windows schedule launcher
changes two path-key expressions to POSIX spelling; all mathematical
conditions and original archive contents are preserved.

\paragraph{Weighted localization and the 34 finite profiles.}
Appendix~\ref{ut:section} gives the mathematical proof of
\eqref{loc:weighted-transfer}, including the raw first step, all LOW and
HIGH terms, the weighted induction and the order of actual-index support
acquisition. The earlier fixed-floor derivations are retained in
\path{companions/localization/analytic/UNIFORM25_CANTELLI_PRODUCTGAP_REVIEW.txt}
and
\path{companions/localization/analytic/UNIFORM25_BETA0_SOURCE_AND_COMPOSITION_REVIEW.txt}.
The inward-plane scaling is derived in
\path{companions/localization/original/new_audit/SCALED_INWARD_PHYSICAL_DERIVATION_V1.txt}.
Their continuous foundation is the resource-enclosure and simultaneous-contract
argument identified above; actual-index support acquisitions are still
separate inputs.

The initial B0/B1 radii now follow directly from
Proposition~\ref{gc:initial-support}. Their original actual-index
radius computations are no longer premises of this first step;
later support refinements still use their own incoming data.
The Peano-kernel and rational-ray calculations are also retained in
\path{gaussian_contact_peano_20261028.tex} and
\path{gaussian_contact_peano_20261028.py} in the finite analytic directory.

The full finite comparison text is
\path{companions/finite-steps/source/analytic/finite_gap_shared_stoploss_20261026.tex}.
Theorems \path{global-c26-complete}, \path{global-c27-complete},
\path{global-c28-complete}, \path{global-c29-complete}, and
\path{global-c30-complete} give the five actual-index compositions.
Their proofs identify the constituent profile lemmas.
At $26$, these include \path{p34-p35-count8-complete},
\path{p35-count9-second-moment-complete},
\path{p36-p37-count10-second-moment-complete}, and the complete-panel
corollaries for $36$--$39$.
At $27$, the profile conclusions are
\path{n27-p34-complete}, \path{n27-p35-complete},
\path{n27-p36-complete}, \path{n27-central-complete}, and
\path{n27-p39-complete}, \path{n27-p40-complete};
the central corollary covers both $37$ and $38$.
The corresponding $28$ statements are
\path{n28-p34-p35-complete}, \path{n28-p36-complete},
\path{n28-p37-complete}, \path{n28-p38-complete}, and
\path{n28-p39-p40-complete}.
The last two groups are the lemmas
\path{n29-seven-panels-complete} and
\path{n30-seven-panels-complete}.
The lower-moment replacement in Lemma~\ref{fc:low40} is the
separate program \path{probe_b0_shared_reference_20261027.py}
in that analytic source directory, with the exact output
\path{result_b0_n26_p40_complete_frozen_20261027.json}.
Its \texttt{--complete} entry recomputes this whole conditional
profile without reading a stored numerical result.
Lemma~\ref{fc:low39} uses the separate
\path{probe_b0_p39_shared_20261027.py} and
\path{b0_allocation_interval_clip_20261027.py}, with output
\path{result_b0_n26_p39_complete_frozen_20261027.json} in the
same analytic directory.
The signed-price extension in Lemma~\ref{fc:low41} is
\path{probe_b0_p41_shared_20261027.py}, with
\path{b0_two_endpoint_budget_20261027.py} and
\path{b0_signed_allocation_interval_20261027.py}; its result is
\path{result_b0_n26_p41_complete_frozen_20261027.json}.
Corollary~\ref{fc:n27-low41} reuses these helpers through the
explicit \texttt{PLAN} in
\path{probe_b0_n27_p41_shared_20261027.py}; the actual-$27$ output is
\path{result_b0_n27_p41_complete_fresh_20261027.json}.
The actual-$28$ plan is
\path{probe_b0_n28_p41_shared_20261028.py}, with output
\path{result_b0_n28_p41_complete_fresh_20261028.json}.
The crossing application at $29$ is the \texttt{PLAN} in
\path{probe_b0_n29_p41_shared_20261028.py}, using
\path{b0_variance_endpoint_branches_20261028.py}; its output is
\path{result_b0_n29_p41_complete_fresh_20261028.json}.
The actual-$30$ plan is
\path{probe_b0_n30_p41_shared_20261028.py}, with final output
\path{result_b0_n30_p41_complete_frozen_20261028.json}.
Its earlier failed complete candidate is retained in the research
sources and supplies no completed profile bound.
The finite lower-envelope implementation in
\path{probe_classical_order_compression_20261028.py} replaces optimization
over mixture weights. Its exact algebraic replay of all $25$ nonempty
N27/P35 mixed regions agrees with the old optimized bound. All other
old fixed-mixture bounds transfer by the pointwise domination following
Lemma~\ref{pp:mixture}; no numerical regeneration is required for this
logical transfer.
The acquisition, lower-moment, ordinary-forest and unrestricted
Fourier inputs and returned objects are retained separately in
\path{companions/finite-other-support/}. Its
\path{read_finite_other_support.py} provides the original mathematical
readers with temporary source materialization. The expensive complete
lower-moment reader is an explicit option. Partial historical count
runs retain their original domains; the analytic profile comparisons
supply the stated complementary regions.

The companion text identifies the continuous inequalities and compositions.
The frozen coefficient data, closed subdivision trees, exact arithmetic
programs and archived fresh calculations are retained in the collections
above. Their mathematical readers relate returned objects to the complete
parameter covers and consumed contracts. Fresh numerical regeneration
is a separate operation using the retained producer entrypoints.
